% INTRODUÇÃO

\chapter{INTRODUÇÃO}
\label{chap:introducao}

Algoritmos e Estruturas de Dados constituem uma área fundamental na formação em Computação, pois fornecem a base para organizar, processar e recuperar informações de maneira eficiente \cite{leblanc2015se2014,cormen2022introduction}. Entre os conteúdos dessa disciplina, Heap e Tabelas Hash são estruturas amplamente utilizadas e exigem articulação entre teoria, análise de complexidade e implementação prática \cite{cormen2022introduction}.

Nos últimos anos, ferramentas de inteligência artificial generativa, como o ChatGPT, passaram a integrar o cotidiano dos estudantes de programação, auxiliando na explicação de conceitos, na geração de exemplos e na sugestão de soluções \cite{baidoo2023education,rahman2023education,garcia2025programming}. Ao mesmo tempo, seu uso suscita questionamentos sobre autonomia, pensamento crítico e possível dependência de respostas prontas \cite{dwivedi2023chatgpt,delikoura2025superficial,qin2023general}, especialmente em disciplinas que exigem raciocínio lógico e a implementação de algoritmos.

Diante desse cenário, o problema investigado neste trabalho foi formulado por meio da seguinte questão de pesquisa: \textit{o uso do ChatGPT durante atividades práticas de Algoritmos e Estruturas de Dados altera o desempenho dos estudantes, medido pela pontuação nos questionários sobre Heap e Tabelas Hash, quando comparadas as condições com e sem a ferramenta?} Essa questão insere-se em investigações recentes sobre o impacto do ChatGPT em programação e em disciplinas de estruturas de dados \cite{jamie2024dsa,kosar2024chatgpt}.

Este Trabalho de Conclusão de Curso foi motivado por um artigo recentemente aprovado no Congresso Brasileiro de Engenharia de Software (CBSoft) \cite{neto2026chatgpt}, que investiga o impacto do ChatGPT na aprendizagem de Algoritmos e Estruturas de Dados. A presente pesquisa dá continuidade a essa linha de investigação, com foco nos conteúdos de Heap e Tabelas Hash.

\section{Objetivos}

\subsection{Objetivo geral}

O objetivo geral deste trabalho é analisar se o uso do ChatGPT durante atividades práticas influencia o desempenho dos estudantes, medido pela pontuação nos questionários sobre Heap e Tabelas Hash.

\subsection{Pergunta de pesquisa}

A partir do objetivo geral, formula-se a seguinte pergunta de pesquisa, alinhada ao problema apresentado na seção anterior: \textit{o uso do ChatGPT durante atividades práticas de Algoritmos e Estruturas de Dados altera o desempenho dos estudantes, medido pela pontuação nos questionários sobre Heap e Tabelas Hash, quando comparadas as condições com e sem a ferramenta?}

\subsection{Objetivos específicos}

Os objetivos específicos são:

\begin{itemize}
    \item caracterizar a participação e o perfil dos estudantes quanto ao uso prévio de IA generativa;
    \item comparar o desempenho nos questionários entre as condições com e sem uso do ChatGPT;
    \item discutir os resultados à luz da literatura e das limitações do desenho experimental adotado.
\end{itemize}

Para atingir esses objetivos, adota-se um desenho experimental com grupos cruzados, no qual cada grupo utiliza a ferramenta em um tópico e não a utiliza no outro. O desempenho é operacionalizado como a pontuação obtida em questionários objetivos aplicados ao final de cada atividade prática, sem apoio do ChatGPT.

